% !TEX program = xelatex
% This is my resume
% Chinese translation
% by Kenji Mouri

\documentclass{resume}

\usepackage{lastpage}
\usepackage{fancyhdr}
\usepackage{linespacing_fix} % disable extra space before next section
\usepackage[fallback]{xeCJK}

\linespread{1.11} 

\begin{document}
\renewcommand\headrulewidth{0pt}

\name{Kenji Mouri}

\basicInfo{
  \email{Mouri\_Naruto@Outlook.com} \textperiodcentered\
  \phone{182-6175-7560} \textperiodcentered\
  \github[MouriNaruto]{https://github.com/MouriNaruto} \textperiodcentered\
}

\section{Education}

\datedsubsection{\textbf{Changshu Institute of Technology}, China}{2016.09 -- 2020.07}
  Major: Automotive Service Engineering, Graduation Date: 07/2020.

\section{Work Experience}
  
\datedsubsection{\textbf{The 301 Lab, Changshu Institute of Technology}}{2018.03 -- 2020.09}
\role {Vehicle-to-everything and Self-driving}{Core Contributor, Future Cleaner Technology Software Engineer}
\begin{itemize}
  \item 独立设计适用于智小蜂无人清扫车的车联网平台，使得适配新新车型的工作量大幅缩减。
  \item 首次在团队提出使用支持响应式布局的现代 Web 技术构建跨设备跨平台的人机交互界面的理念并付诸实践，在此期间获得了 ASP.NET Core 和现代 Web 技术的开发经验，并实现了基于 Xamarin 的 Android 客户端。
  \item 参与第十六届“挑战杯”，获取江苏省大学生课外学术科技作品竞赛二等奖。
  \item 协助开发实验室物联网平台，使其支持西门子 S7-1200 PLC 并提供二次开发 SDK，对大屏展示界面进行了改进。
\end{itemize}

\datedsubsection{\textbf{Chuyu Team}}{2014.07 -- Present}
\role {Windows Application Development}{Core Contributor}
\begin{itemize}
  \item 协助开发 \href{http://www.chuyu.me/zh-Hans/index.html}{DISM++} 项目，引入 NCleaner 清理增强组件、映像热备份和多语言支持。
  \item 独立开发 \href{https://github.com/Chuyu-Team/MINT}{MINT} 项目，基于 PHNT 项目制作的 Windows 未文档化的用户模式 API 定义。
  \item 独立开发 \href{https://github.com/Chuyu-Team/libkcrt}{libkcrt} 项目，现代 MSVC 工具链下连接到 ntdll.dll 中 CRT 函数的解决方案。
  \item 协助开发 \href{https://github.com/Chuyu-Team/VC-LTL}{VC-LTL} 的 Visual Studio 工具集成并引入 \href{https://github.com/Chuyu-Team/vc-ltl-samples}{VC-LTL 项目的使用示例}。
  \item 协助改善 \href{https://github.com/Chuyu-Team/YY-Thunks}{YY-Thunks (鸭船)} 项目。
\end{itemize}

\section{Personal Projects}

\datedsubsection{\textbf{NSudo}}{\url{https://github.com/M2Team/NSudo}}
一套系统管理工具集，NSudo Launcher 用于在 Windows 下以特定权限和设置运行应用，NSudo Devil Mode 是一个为开发者提供在管理员权限下绕过 Windows 的文件和注册表的访问检查的解决方案。

\datedsubsection{\textbf{Nagisa}}{\url{https://github.com/Project-Nagisa/Nagisa}}
一个基于 UWP 平台的文件传输工具，支持后台传输，还处于早期开发阶段。

\datedsubsection{\textbf{Mile}}{\url{https://github.com/MouriNaruto/Mile}}
个人项目使用的公共库，使用 C++17 的编译器语法特性，但提供 C ABI 兼容。

\section{Skills}
\begin{itemize}

  \item \textbf{Program Language}:
    \textbf{multilingual}(not limited to any specific language), usually use C/C++, C\# and Python.

  \item \textbf{Windows Application Development}:
    \textbf{6 years} of experience, familiar with Windows SDK (Win32, COM, ATL, WTL, WRL, C++/CX and C++/WinRT).

  \item \textbf{.NET Application Development}:
    \textbf{2 years} of experience, familiar with WPF, WinForms, ASP.NET Core and Xamarin.
  
  \item \textbf{Linux Application Development}:
    \textbf{2 years} of experience, familiar with ROS (Kinetic and Melodic) and Qt.

  \item \textbf{Development Tool}:
    \textbf{can adapt to any mainstream development tools under Windows and Linux}, usually use Visual Studio and IDA Pro under Windows, or Clion and PyCharm under Linux, also use Git and VSTS.

\end{itemize}

\section{Miscellaneous}
\begin{itemize}

  \item Languages: English - good (passed CET-6),  Chinese - native speaker.

  \item I am interested in writing the most compact implementations by using the least syntaxes and third-party libraries.

  \item Opensource Contributions: \href{https://github.com/pulls?q=is:pr+author:MouriNaruto}{https://github.com/pulls?q=is:pr+author:MouriNaruto} contributed to OpenSSL, Windows File Manager, Covariant Script, FFmpegInteropX and other projects.
  
  \item 在 My Digital Life Forums、远景论坛、吾爱破解论坛、看雪安全论坛以 Mouri\_Naruto 发表过精华或置顶帖。
  
  \item 渴望结识更多的可以进行技术交流的友人。

  \item Another style of my introduction (Chinese version): \href{https://mourinaruto.github.io/about}{https://mourinaruto.github.io/about}。
  
\end{itemize}

\end{document}
